\documentclass[10pt]{article}
\usepackage[utf8]{inputenc}
\usepackage[T1]{fontenc}
\usepackage{amsmath}
\usepackage{amsfonts}
\usepackage{amssymb}
\usepackage[version=4]{mhchem}
\usepackage{stmaryrd}
\usepackage{bbold}

\begin{document}
$$
L 7
$$

\section*{L7: The Chinese Remainder Theorem}
Example:- Find a positive integer having a remainder of 2 when divided by 3, a remainder of 1 when divided by 4 , and a remainder of 3 when divided by 5 .\\
In other words, we want to find a positive integer $x$ such that:

$$
\begin{aligned}
& x \equiv 2(\bmod 3) \\
& x \equiv 1(\bmod 4) \\
& x \equiv 3(\bmod 5)
\end{aligned}
$$

This is a system of linear congruences in one variable.\\
Theorem 2.9 (Chinese Remainder Theorem):\\
Let $m_{1}, m_{2}, \ldots, m_{n}$ be pairwise velatively prime positive integers and let $b_{1}, \ldots, b_{n}$ be any integers. Then the system of linear congruences in one variable given by

$$
\begin{aligned}
& x \equiv b_{1}\left(\bmod m_{1}\right) \\
& x \equiv b_{2}\left(\bmod m_{2}\right)
\end{aligned}
$$

$$
x \equiv b_{n}\left(\bmod m_{n}\right)
$$

has a unique solution modulo $m_{1} m_{2} \ldots m_{n}$.\\
Proof: We first construct a solution. Let

$$
\begin{aligned}
& M:=m_{1} m_{2} \ldots m_{n}, \text { and } \\
& M_{i}:=M_{1} m_{i} \text { for } i=1,2, \ldots, n .
\end{aligned}
$$

$\operatorname{Thus}\left(M_{i}, m_{i}\right)=1$ for each $i$, and hence $m_{i} x_{i} \equiv 1\left(\bmod m_{i}\right)$ has a solution for each $i$ by Corollary 2.8. Consider

$$
x=b_{1} M_{1} x_{1}+b_{2} M_{2} x_{2}+\ldots+b_{n} M_{n} x_{n} .
$$

This is a solution to the desired system since for each $i=1,2, \ldots, n$ we have

$$
\begin{aligned}
x & =b_{1} M_{1} x_{1}+\cdots+b_{i} M_{i} x_{i}+\cdots+b_{n} M_{n} x_{n} \\
& \equiv 0+\cdots+b_{i}+\cdots+0\left(\bmod m_{i}\right) \\
& \equiv b_{i}\left(\bmod m_{i}\right) .
\end{aligned}
$$

It remains to show uniqueness of the solution Modulo M. Let $x^{\prime}$ be another solution to the given system of linear congruences. Then we have that $x^{\prime} \equiv b_{i}\left(\bmod m_{i}\right)$ for all $i$.

Since $x \equiv b_{i}\left(\bmod m_{i}\right)$, we have that $x \equiv x^{\prime}$ (mod $m_{i}$ ) for all $i$, or equivalently $\rightarrow m_{i} \mid x-x^{\prime}$ for all $i$. Thus $M \mid x-x^{\prime}$ since $m_{1}, m_{2}, \ldots, m_{n}$ are pairwise velatively prime. Thus $x \equiv x^{\prime}(\bmod M)$, and the proof is complete.

Example (continued): Recall that we needed to Solve:

$$
\begin{aligned}
& x \equiv 2(\bmod 3) \\
& x \equiv 1(\bmod 4) \\
& x \equiv 3(\bmod 5)
\end{aligned}
$$

Using the notation of Theorem 2.9 we have

$$
\begin{aligned}
& M=3 \cdot 4 \cdot 5=60 \\
& M_{1}=M / M_{1}=60 / 3=20 \\
& M_{2}=M / m_{2}=60 / 4=15 \\
& M_{3}=M / m_{3}=60 / 5=12 .
\end{aligned}
$$

We now solve the linear congruences $M_{i} x_{i} \equiv 1\left(\bmod m_{i}\right)$ for $i=1,2,3$. In general, one would use the Euclidean algor ithm to\\
to solve these congrnences. We solve each one by inspection (apribsi knowing that\\
$M_{i} x_{i} \equiv 1\left(\bmod m_{i}\right)$ has a unique solution by\\
Conollary 2.8). Observe that $M_{1} x_{1} \equiv 1\left(\bmod m_{1}\right)$\\
becomes $20 x_{1} \equiv 1(\bmod 3)$\\
iff $\quad 2 x_{1} \equiv 1(\bmod 3)$\\
iff $\quad x_{1} \equiv 2(\bmod 3)$.\\
Observe that $M_{2} x_{2} \equiv 1\left(\bmod m_{2}\right)$ becomes\\
iff\\
$15 x_{2} \equiv 1(\bmod 4)$\\
$3 x_{2} \equiv 1(\bmod 4)$\\
iff\\
$x_{2} \equiv 3(\bmod 4)$\\
Finally, observe that $M_{3} x_{3}\left(\bmod M_{3}\right)$ becomes\\
iff\\
$12 x_{3} \equiv 1(\bmod 5)$\\
$2 x_{3} \equiv 1(\bmod 5)$\\
ift $\quad x_{3} \equiv 3(\bmod 5)$.\\
Thus we put $x_{1}=2, x_{2}=3, x_{3}=3$. The desired solution is

$$
\begin{aligned}
x & =b_{1} M_{1} x_{1}+b_{2} M_{2} x_{2}+b_{3} M_{3} x_{3} \\
& =(2)(20)(2)+(1)(15)(3)+(3)(12)(3) \\
& =233 \\
& \equiv 53(\bmod 60) .
\end{aligned}
$$

The minimum positive solution to the System is 53.

\section*{Wilson's Theorem}
Lemma 2.10: Let $p$ be a prime and let $a \in \mathbb{Z}$. Then $a$ is its own inverse modulo $p$ if and only if $a \equiv \pm 1(\bmod p)$.\\
Proof: $(\Rightarrow)$ Assume that $a$ is its own inverse modulo $p$. Then $a^{2} \equiv 1(\bmod p)$. Thus $p \mid a^{2}-1$, or equivalently, $p((a-1)(a+1)$. By Lemmal. 14 we have that $p \mid a-1$ or $p \mid a+1$. Thus $a \equiv \pm 1(\bmod p)$, as desired.\\
$(\Leftrightarrow)$ Assume that $a \equiv \pm 1(\bmod p)$. Then $a^{2} \equiv 1(\bmod p)$, and $a$ is its own inverse modulo p, as desired.

Theorem 2.11 (Wilson): Let $p$ be a prime. (6) Then $(p-1)!\equiv-1(\bmod p)$.

\begin{itemize}
  \item John wilson conjectured the result based on numerical evidence, but could not prove it.
  \item The Theorem was first proved by Lagrange in 1771.\\
Example (idea behind the proof): We prove that $(11-1)!\equiv-1(\bmod 11)$. We have
\end{itemize}

$$
10!=1 \cdot 2 \cdot 3 \cdot 4 \cdot 5 \cdot 6 \cdot 7 \cdot 8 \cdot 9 \cdot 10
$$

By Lemma 2.10 , the only integers between 1 and 10 inclusive that are their own inverses modulo II are I and 10. We now take each integer from 2 to 9 inclusive and pair it with its inverse modulo II:

$$
\begin{aligned}
& 2 \leftrightarrow 6 \\
& 3 \leftrightarrow 4 \\
& 5 \leftrightarrow 9 \\
& 7 \leftrightarrow 8
\end{aligned}
$$

The product of each pair is congruent to

I modulo II, from which

$$
\begin{aligned}
10! & =1 \cdot 2 \cdot 3 \cdot 4 \cdot 5 \cdot 6 \cdot 7 \cdot 8 \cdot 9 \cdot 10 \\
& =1 \cdot(2 \cdot 6)(3 \cdot 4)(5 \cdot 9)(7 \cdot 8) \cdot 10 \\
& \equiv 1 \cdot 1 \cdot 1 \cdot 1 \cdot 1 \cdot(-1)(\bmod 11) \\
& \equiv-1(\bmod 11) .
\end{aligned}
$$

as desired.\\
Proof (of Wilson's Theorem): The statement of the Theorem is easily checked for $p=2$ and $p=3$. Assume that $p>3$. Each $a \in \mathbb{Z}$ with $1 \leqslant a \leqslant p-1$ has a unique inverse modulo $p$ by Corollary 2.8. Without loss of generality, we assume that the inverse of such an $a$, say $a^{\prime}$, satisfies $1 \leqslant a^{\prime} \leqslant p-1$. Now, by Lemma 2.10, the only integers that are their own inverses modulo $p$ are those integers congruent to $\pm 1(\bmod p)$, so the only integers a with $1 \leq a \leq p-1$ that are their own inverses modulo $p$ are 1 and $p-I$. Therefore\\
for each integer a with $2 \leqslant a \leqslant p-2$, there exists a unique different integer $a^{\prime}$ with $2 \leqslant a^{\prime} \leqslant p-2$ such that $a a^{\prime} \equiv 1(\bmod p)$. Thus we can group the numbers $2,3, \ldots, p-2$, into pairs $a, a^{\prime}$, such that

$$
a a^{\prime} \equiv 1(\bmod p) . \quad(*)
$$

The product of the left sides of the $\frac{p-3}{2}$ congruences $(*)$ is $(p-2)$ ! after rearrangement. The right sides multiply to 1. Thus

$$
(p-2)!\equiv 1(\bmod p) \cdot(* *)
$$

Multiplying both sides of (*) by $p-1$ we obtain

$$
(p-2)!(p-1) \equiv p-1(\bmod p),
$$

and thus

$$
(p-1)!\equiv-1(\bmod p),
$$

as desired.

\begin{itemize}
  \item The converse of Wilson's Theorem is also true.
\end{itemize}

Proposition 2.12: Let $n \in \mathbb{Z}$ with $n>1$. If $(n-1)!\equiv-1(\bmod n)$, then $n$ is a prime. Proof: Let $n=a b$ where $a, b \in \mathbb{Z}$ with $1 \leqslant a<n$. It suffices to show that $a=1$, so that any such factorisation of $n$ is trivial (consisting of 1 and $n$ only). Since $1 \leqslant a<n$, we have that a $\mid(n-1)!$. Now the hypothesis $(n-1)!\equiv-1(\bmod n)$ implies that $n \mid(n-1)!+1$. Since $a \mid n$ we have that a $\mid(n-1)!+ \pm$ by Proposition 1.1. Thus a| $(n-1)!+1-(n-1)!=1$ by Proposition 1.2 and thus $a \mid 1$. So $a=1$, as required.

Remark: •Proposition 2.12 gives a primality test! Although, it is very ineff.cet. inefficient i.e. Computation of $(n-1)$ ! requires at least $n-3$ multiplications modulo $n$ )

\begin{itemize}
  \item wilson's name is also cessociated to a\\
class of primes.
\end{itemize}

Definition: A prime $p$ is said to be a wilson prime if $(p-1)!\equiv-1\left(\bmod p^{2}\right)$.

\begin{itemize}
  \item Only three wilson primes are known:\\
$5,13,563$. It is not known whether\\
there are finitely or infinitely many wilson primes.
\end{itemize}

Fermat's Little Theorem\\
Theorem 2.13 (Fermat's Little Theorem): Let $p$ be a prime and $a \in \mathbb{Z}$. If $p+a$, then

$$
a^{p-1} \equiv 1(\bmod p) .
$$

Proof: Consider the $p-1$ integers given by $a, 2 a, \ldots,(p-1) a$. Note that $p+i a$ for $i=1,2, \ldots, p-1$ (if plia for some $i$, then pla or pli by Lemma 1.14, which is impossible). Note also that no two of the given $p-1$ integers are congruent modulo P. Why? Since $p+a$, the inverse of $a_{\text {, }}$\\
say $a^{\prime}$, exists by Corollary 2.8 , so if (11)\\
$i a \equiv j a(\bmod p)$ with $i \neq j$, then\\
i $a a^{\prime} \equiv j a a^{\prime}(\bmod p)$, from which $i \equiv j(\bmod p)$,\\
which is impossible. So the least nonnegative resides modulo $p$ of $a, 2 a, \ldots$,\\
$(p-1) a$, taken in some order must be $1,2, \ldots, p-1$. Then,\\
$(a)(2 a) \cdots((p-1) a) \equiv 1 \cdot 2 \cdots(p-1)(\bmod p)$, or equivalently,

$$
a^{p-1}(p-1)!\equiv(p-1)!(\bmod p) .
$$

By wilson's Theorem, the congruence becomes

$$
-a^{p-1} \equiv-1(\bmod p),
$$

and so $a^{p-1} \equiv 1(\bmod p)$.\\
Remark: Wilson's Theorem is not needed here.\\
why?


\end{document}